\documentclass[11pt,a4paper]{article}

\usepackage[UTF8]{ctex}
\usepackage{amsmath,amssymb}
\usepackage{geometry}
\usepackage{booktabs}
\usepackage{xcolor}
\geometry{margin=2.5cm}

\title{小车控制器架构设计思路}
\author{}
\date{\today}

\begin{document}
\maketitle

\section{三层架构: NN $\to$ 虚拟车 $\to$ 跟踪}

\subsection{为什么不能 NN 直接输出舵机/油门}

端到端控制的根本问题: \textbf{NN 的 20Hz 输出是阶梯状的}。

\begin{center}
\begin{tabular}{c|c|c}
\toprule
& NN 直接输出 & NN $\to$ vst $\to$ 跟踪 \\
\midrule
控制量 & 舵角、油门 (位置量) & 加速度、转向速率 (高阶量) \\
NN 输出 & 20Hz 跳变 & 20Hz 跳变 \\
执行器指令 & 20Hz 跳变 & \textbf{200Hz 平滑} \\
路径连续性 & 阶梯状 & 连续 \\
\bottomrule
\end{tabular}
\end{center}

如果 NN 直接输出舵角，20Hz 的零阶保持意味着每 50ms 舵机突然跳到新位置。舵机有能力 200Hz 平滑过渡，但 NN 不给它这个机会。

\subsection{解决方案: NN 输出高阶控制量}

NN 不输出位置 ($\delta$, $v$)，而是输出导数 ($\omega$, $a$)：

\begin{equation}
\begin{bmatrix} \omega \\ a \end{bmatrix}_{\text{20Hz}} = \text{NN}(v, \delta, \mathbf{g}_1, \mathbf{g}_2, \mathbf{g}_3)
\end{equation}

虚拟车在 200Hz 下积分:
\begin{align}
\delta_{\text{vst}}(t) &= \int \omega_{\text{NN}} \,dt \\
v_{\text{vst}}(t)     &= \int a_{\text{NN}} \,dt \\
\theta_{\text{vst}}(t) &= \int \frac{v_{\text{vst}}}{L}\tan\delta_{\text{vst}} \,dt \\
x_{\text{vst}}(t)     &= \int v_{\text{vst}}\cos\theta_{\text{vst}} \,dt \\
y_{\text{vst}}(t)     &= \int v_{\text{vst}}\sin\theta_{\text{vst}} \,dt
\end{align}

\textbf{NN 的 20Hz 阶梯被积分器平滑成 200Hz 的连续轨迹。} 这就是选择高阶控制量的核心原因。

\subsection{为什么需要虚拟车}

虚拟车的三个角色:

\begin{enumerate}
\item \textbf{平滑器}: 把 NN 的阶梯输出积分成连续路径
\item \textbf{参考生成器}: 产生 $(x,y,\theta,v,\delta)_{\text{vst}}$ 五个参考状态
\item \textbf{误差坐标系}: 在 vst 局部坐标系下投影误差，解耦纵向/横向
\end{enumerate}

没有虚拟车的替代方案:
\begin{itemize}
\item NN 直接输出轨迹点 $\to$ 需要二次插值 $\to$ 增加延迟和复杂度
\item NN 输出高阶控制量直接给执行器 $\to$ 无反馈修正，开环
\end{itemize}

虚拟车是最小化的中间层——只做运动学积分，不引入额外参数。

\section{前馈 + 反馈}

\subsection{分工}

\begin{center}
\begin{tabular}{c|c|c}
\toprule
& 前馈 ($\omega_{\text{ff}}$, $a_{\text{ff}}$) & 反馈 (LQR, LADRC) \\
\midrule
来源 & NN 输出 & 传感器误差 \\
作用 & 预期轨迹跟踪 & 扰动抑制 \\
延迟 & 0 (零阶保持) & 5ms (传感器 + 计算) \\
频带 & 低频 (0--20Hz) & 中高频 (0--100Hz) \\
失效时 & 开环跑偏 & 仍能稳定 \\
\bottomrule
\end{tabular}
\end{center}

前馈承担 80\% 的控制力——NN 根据航点规划了全局最优路径，前馈就是这条路径的执行指令。反馈只修正剩余的 $\sim$5--10cm 误差。

\subsection{为什么不全用反馈}

纯反馈 (如纯 LQR 跟踪航点) 的问题:
\begin{itemize}
\item 反馈只看当前误差，不知道前方弯道——入弯晚了才反应
\item 增益必须覆盖最坏情况 (急弯 + 高速)，日常场景下过于激进
\item 没有前馈的 "预判"，在 U 型弯等大曲率路段会明显滞后
\end{itemize}

\subsection{为什么不全用前馈}

纯前馈 (NN 输出直接给执行器) 的问题:
\begin{itemize}
\item NN 只有 20Hz，高频扰动 (地面不平、小石子) 来不及响应
\item 模型误差 (轮胎滑移、陀螺零偏) 会累积，无纠正机制
\item 传感器数据不进回路，开环精度依赖模型精度
\end{itemize}

前馈 + 反馈本质上是 \textbf{低频规划 + 高频纠偏}。NN 负责 "往哪走"，LQR/LADRC 负责 "走准它"。

\section{为什么用 200Hz 跟踪层}

\begin{center}
\begin{tabular}{c|c}
\toprule
频率 & 理由 \\
\midrule
20Hz  & NN 推理 ($\sim$10ms) + 通信 ($\sim$10ms) $\approx$ 20Hz 上限 \\
200Hz & 舵机带宽 $\sim$200Hz，5ms 控制周期足够跟踪瞬态 \\
\bottomrule
\end{tabular}
\end{center}

跟踪层跑在 ISR 里，200Hz 不动摇。NN 跑在 main 线程，20Hz 够用。两层之间通过 `PlannerAction` 结构体零阶保持传递——NN 更新了就接新值，没更新就保持旧值。

ZOH 的合理性: NN 输出的 $\omega$ 和 $a$ 是\textbf{导数}，ZOH 保持意味着虚拟车以恒定角速度和加速度运动——这是物理上自然的短时假设 (车辆不能在 50ms 内突变加速度)。

\section{误差投影: 为什么用 vst 局部坐标系}

\begin{equation}
\begin{bmatrix} e_x \\ e_y \end{bmatrix} =
\begin{bmatrix}
\cos\theta_{\text{vst}} & \sin\theta_{\text{vst}} \\
-\sin\theta_{\text{vst}} & \cos\theta_{\text{vst}}
\end{bmatrix}
\begin{bmatrix}
x_{\text{vst}} - x_{\text{car}} \\
y_{\text{vst}} - y_{\text{car}}
\end{bmatrix}
\end{equation}

\begin{center}
\begin{tabular}{c|c}
\toprule
世界坐标系误差 & vst 局部坐标系误差 \\
\midrule
$e_x$, $e_y$ 耦合 & $e_x$ = 纵向误差, $e_y$ = 横向误差, 解耦 \\
转弯时 $e_x$ 和 $e_y$ 都大 & 转弯时参考坐标系跟着转, 误差自然分解 \\
LQR 需要解耦模型 & LQR 模型假设纵向/横向独立 \\
\bottomrule
\end{tabular}
\end{center}

\section{为什么用速度自适应 LQR}

LQR 增益表在 11 个速度断点上离线求解，运行时间距查找 + 双线性插值:

\begin{equation}
\mathbf{K}(v) = \text{LQR}\big(\mathbf{A}(v), \mathbf{B}, \mathbf{Q}, \mathbf{R}\big)
\end{equation}

$\mathbf{A}$ 矩阵的 $(1,2)$ 元素是 $v$——速度变了，系统动态变了。固定增益 LQR 在低速过阻尼、高速欠阻尼。速度自适应增益在全速域保持一致的闭环特性。

\section{为什么用 LADRC 而非 PID}

纵向的三个主要扰动源:

\begin{enumerate}
\item 坡度: 上坡时有效加速度减少 $g\sin\phi$
\item 摩擦: 地面材质变化 (瓷砖 vs 地毯)
\item 电池电压: 电压跌落导致电机力矩下降
\end{enumerate}

PID 通过积分项逐步补偿，响应慢且容易积分饱和。LADRC 的 LESO 将\textbf{所有扰动汇总为单一量 $f$}，在观测器带宽 (15 rad/s, $\sim$40ms) 内实时估计并前馈补偿。

降阶 LESO 的选择: 速度 $v$ 已被编码器精确测量，不必重复观测。仅观测扰动 $f$ 即可，降低观测器阶数、提高收敛速度。

\section{总结: 架构全景}

\[
\begin{array}{c}
\fbox{NN 规划器 (20Hz)} \\
\downarrow \quad [\omega_{\text{ff}},\; a_{\text{ff}}] \\
\fbox{虚拟车积分器 (200Hz)} \\
\downarrow \quad [x,y,\theta,v,\delta]_{\text{vst}} \\
\fbox{误差投影 (vst 坐标系)} \\
\downarrow \quad [e_x, e_y, e_\theta, \dot{e}_\theta, e_\delta] \\
\begin{array}{|c|c|}
\hline
\fbox{LQR 横向} & \fbox{LADRC 纵向} \\
\text{6-term 增益插值} & \text{降阶 LESO + PD} \\
\omega_{\text{cmd}} = \omega_{\text{ff}} + \mathbf{Kx} & u = (k_p e_x + \cdots - \hat{f})/b_0 \\
\hline
\end{array} \\
\downarrow \quad [\delta_{\text{cmd}},\; u_L, u_R] \\
\fbox{执行器 (舵机 + 电机)} \\
\downarrow \\
\fbox{真实车辆} \\
\downarrow \quad [\text{IMU}, \text{编码器}] \\
\uparrow \quad \text{反馈至误差投影}
\end{array}
\]

\textbf{一句总结: NN 说往哪走 (前馈)，物理模型连续化路径 (虚拟车)，传感器说实际在哪 (反馈)，LQR/LADRC 弥合差距 (纠偏)。}

\end{document}
